\section{Cleanup of the unresolved boundary}
\label{sec:cleanup}

The early phase leaves fewer than $b$ bad stubs.  The separation
construction in Lemma~\ref{lemma:structure} resolves them in $b$ steps,
from $l_1=l_0+2b$ unresolved stubs to $l_0$.  We now prove that this
cleanup also preserves the analytic estimates needed for boundary
completion.

Fix an order for the initially bad stubs and use them as first endpoints.
After they are resolved, take the first still-surviving stub in a fixed
order.  At each step the first endpoint $u$ is prescribed and the second
endpoint $v$ is uniform among
the eligible surviving initially clean stubs.  Different steps use
distinct first stubs.  The eligible set omits only $O(b)$ currently
unresolved stubs, and every resulting pair satisfies the distance and
simplicity requirements of Lemma~\ref{lemma:structure}.

The prescribed endpoint changes the analytic argument: a two-endpoint
average is no longer available.  Instead, sums along the distinct first
endpoints will be bounded by path sums in the graph at the end of the
already chosen history.  This gives the cumulative estimates needed by
the account rule.

Write $f_b=b/l_0$.  The parameter choices give
\begin{equation}
\begin{gathered}
 f_b=o(1),\qquad b^{3/2}/l_0=o(1),\qquad
 f_b\sqrt{l_0}\asymp n^{.27a},\\
 f_b n^\eta q^{L/2}=o(1),\qquad
 f_b n^{v_*}/\alpha_*=o(1),\qquad
 \alpha_*\sqrt{l_0}\asymp n^{.575a}.
\end{gathered}
\label{eq:original-35}
\end{equation}
Whenever a small loss $n^\eta$ occurs in this section, its exponent may
be prescribed in advance: the structural and preceding-phase losses
are first chosen sufficiently smaller.  There are only finitely many
uses of this convention.

For an output label $i$, assign a level $\nu(i)$: level $2$ for a global
vertex coordinate, level $1$ for an initially bad stub, and level $0$ for
an initially clean stub.  A stub label is used only while it survives.
The total level of an entry $E_{ij}$ is $\nu(i)+\nu(j)$.  Set
\[
 B_c=\frac{n^{.35a}}{\sqrt{l_0}},\qquad
 B_k=B_c n^{4k\rho}\quad(0\leq k\leq4),\qquad B_m=2B_4.
\]
Choose $\rho>0$ sufficiently small that
$B_m\leq n^{.37a}/\sqrt{l_0}$ for all sufficiently large $n$.
The bootstrap bound for an entry of total level $k$ is $2B_k$.
The small increase in scale between successive levels will pay for
transfers from a global coordinate to a stub, or from a bad stub to a
clean stub, in the inverse update.
All primary and auxiliary values of $h$, and both signs, remain in use.

\begin{samepage}
\begin{proposition}[Cleanup estimates]\label{prop:cleanup}
Assume the early-phase estimates at $l_1$, and maintain the structural
estimators during cleanup.  With the simultaneous discounted-account
choice rule, the cleanup interval can be completed so that, at every
chosen state, the precision matrices remain positive definite,
\eqref{eq:original-8} holds, and
\[
 T_S=\sum_{x\in S}e_x(x)^2\leq n^\eta,\qquad
 |E_{ij}|\leq2B_k.
\]
The second inequality holds for every surviving pair $(i,j)$ of total
level $k=\nu(i)+\nu(j)$.
At primary points the lower estimate for $\mathcal T$ persists in the
form
\[
 \mathcal T\geq n^{-\theta}(n^{-1}+\alpha/l)^{-1}.
\]
In particular the consequences \eqref{eq:original-10} and
\eqref{eq:original-11}, and a polynomial-loss form of
\eqref{eq:original-13}, hold at $l_0$.
\end{proposition}
\end{samepage}

The proof has three parts.  First we establish pathwise summation bounds
and preserve the scalar traces.  Next we classify the inverse-update
terms and control the surviving diagonal-square sum $T_S$.  Finally we
bound individual entry increments and extract all estimates
simultaneously.  Each one-step estimate assumes only the current
bootstrap bounds.
Throughout the proof, $l\asymp l_0$; in cumulative upper bounds it is
therefore harmless to replace $l$ by $l_0$.

\subsection{Majorants along a chosen history}

We retain the absolute path and endpoint majorants
$A_i(x),j_x,b_{xy}$ of Lemma~\ref{lemma:path-majorants}.  Define
\[
 r_i(x)=A_i(x)+|e_i(x)|,\qquad Q_a=n^\eta q^{L/2}.
\]
On current stub coordinates the arrays $r_x(y)$ may be taken
symmetric.  Their row and column squared sums are bounded, since this
is true for $A$, $P$, and $W$.  Consequently the symmetric nonnegative
matrix $(r_x(y)^2)_{x,y\in S}$ has bounded operator norm: its row sums
are bounded, and the same is true of its column sums.

\begin{lemma}[Pathwise summation]\label{lem:cleanup-pathwise}
Fix any chosen cleanup history up to its current state, and fix a
label $i$ while it is alive.  With $u_t$ the prescribed first endpoint
of step $t$, one has
\begin{equation}
 \sum_t A_i^t(u_t)^2\leq n^\eta,\qquad
 \sum_t (j_{u_t}^t)^2\leq n^\eta.
\label{eq:original-36}
\end{equation}
The first estimate also holds with $A_i^t(u_t)$ replaced by
$r_i^t(u_t)$.  For every fixed initial stub $v$, whether or not it has
already been resolved, the actual-path part of $b_{u_t v}^t$ satisfies
the same squared-sum estimate over $t$.
\end{lemma}

\begin{proof}
Let $G_*$ be the actual graph at the last state of the history segment.
Every actual path present earlier is present in $G_*$.  The first
endpoints are distinct initial stubs, and at most $d$ such stubs lie
at one vertex.  Thus it suffices to bound an absolute weighted path
row to all vertices of $G_*$.  In particular, the target stub need not
remain unresolved in $G_*$.  At a fixed length $k<L$, the total
number of paths is at most $Cq^k$ and their fixed-endpoint
multiplicity is bounded.  Since $h^{2k}\leq q^{-k}$, the squared sum
of the weighted row at that length is bounded.  Cauchy--Schwarz over
the $O(\log n)$ lengths gives a polylogarithmic bound for the total
row.  This argument also applies with a fixed stub vertex as the
starting point and the distinct $u_t$ as targets.

The nonempty-return squared-sum estimate follows in the same way
from the global return-square diagram bound in $G_*$.  Each uniform
or centering contribution is at most
$n^{\eta_0}q^{L/2}/l_0$, with a preliminary loss $\eta_0$.
Even summing its square over all $l_1$ initial stubs gives
\[
 O\!\left(n^{2\eta_0}q^L/l_0\right)=o(1).
\]
Finally,
$\sum_t|e_i^t(u_t)|^2\leq bB_m^2=o(1)$ under the entry bootstrap.
Choosing $\eta_0$ small enough proves all the assertions.
\end{proof}

In particular, Cauchy--Schwarz gives
\[
 \sum_t r_i^t(u_t)\leq n^\eta\sqrt b,\qquad
 \sum_t j_{u_t}^t\leq n^\eta\sqrt b,\qquad
 \sum_t r_i^t(u_t)j_{u_t}^t\leq n^\eta.
\]
The same bounds apply to every prefix ending at the current chosen
state.  They will control the accumulated discounts of the entry
accounts, although the distribution at each step averages only the
second endpoint.  No future completion is used.

\subsection{Trace errors and the lower trace bound}

\begin{lemma}[Persistence of scalar estimates]
\label{lem:cleanup-traces}
At a current cleanup state, assume positive definiteness and
\eqref{eq:original-8} at every tracked point for both signs, together
with the level-dependent entry and $T_S$ bootstrap bounds.
The trace-square accounts at all tracked points and the
$\mathcal T$-deviation accounts at the primaries have bounded discount
loss at scales strictly inside the corresponding bounds of
Proposition~\ref{prop:cleanup}.
\end{lemma}

\begin{proof}
We keep the trace errors $L_1=F_1-(l-1)$ and $L_0=F_0-(n-1)$
at the intermediate scales
\[
 S_1=n^{.085a}/\alpha_*,\qquad S_0=n^{.115a}/\alpha_*^2.
\]
At each primary, we reinitialize $\mathcal T$ with its deviation from
its value at $l_1$, rather than continue the reciprocal account used in the early
phase.  As in the trace calculation there, $K,Q,J$ denote the current
pair blocks of $K_{\rm all},Q_{\rm all},\Gamma Q_{\rm all}$.

\emph{The stub trace $L_1$.}
In the expansion \eqref{eq:original-15}, the signed conditional mean of all terms
other than the ideal term $-\operatorname{Tr}H_0K$ is bounded in
absolute value by
\[
 C\left(j_u+B_m+\frac{n^\eta}{\sqrt l}+\frac1l\right).
\]
Here the bounds \eqref{eq:original-14} control the endpoint factors,
and averaging the random endpoint uses the row cross norms and
$T_S$.  Its cumulative bound is
\[
 C\left(n^\eta\sqrt b+bB_m+
             \frac{n^\eta b}{\sqrt{l_0}}+f_b\right).
\]
For the pure-swap cross term, the full partner sum cancels by
centering when the artificial partner $v=u$ is included.  The full
coefficient $H_0$ has a finite-size remainder.  More precisely,
since $H_0=-\sigma h(F+J_2/(l-2))$, uniform averaging over the
distinct partners gives
\[
 \mathbb E_{v\ne u}[-\operatorname{Tr}(H_0K)]
   =\sigma h\frac{\operatorname{Tr}K_{\rm all}-lK_{uu}}
                       {(l-1)(l-2)}.
\]
Here $\sum_v(K_{\rm all})_{uv}=0$ was used.  The remainder is
$O(1/l)$ because all diagonal entries of $K_{\rm all}$ are bounded.
For the actual eligible set $I$, its complement has $O(b)$ stubs, and
\[
 \|(K_{\rm all})_{u,*}\|_2\leq Cn^{v_*}/\alpha.
\]
Indeed $K_{uu}\leq C$, and the operator bound for the relevant
positive Gram matrix controls its row norm.  The centered cross sum
over $I$ is the negative omitted-row sum.  Cauchy--Schwarz bounds
this sum, while the non-pure diagonal contribution is $O(1/l)$
directly.  Thus the absolute signed mean is at most
\[
 C\left(\frac{\sqrt b\,n^{v_*}}{\alpha l}+\frac1l\right).
\]
The sum of conditional second moments is at most
\[
 C\left(b+f_b\frac{n^{2v_*}}{\alpha^2}\right).
\]
All these mean bounds divided by $S_1$, and all second-moment
bounds divided by $S_1^2$, tend to zero by a fixed power of $n$.

\emph{The global trace $L_0$.}
For $L_0$, the nonideal mean bound above acquires a factor $C/\alpha$.
The ideal $Q$ row has norm at most $Cn^{v_*}/\alpha^2$, by
\eqref{eq:original-12} and $Q_{uu}\leq C/\alpha$.
The same fixed-endpoint calculation has finite-size remainder
$C/(\alpha l)$, and omitted-row contribution
$C\sqrt b\,n^{v_*}/(\alpha^2l)$.
The corresponding cumulative second-moment bound is
\[
 C\left(\frac b{\alpha^2}+
              f_b\frac{n^{2v_*}}{\alpha^4}\right).
\]
Division by $S_0$ and $S_0^2$ gives the same power savings for
the mean and second-moment budgets.  For extraction, if a scalar
increment has signed mean bounded by $\mu_t$ and second moment
bounded by $\nu_t$, the account
$1+L^2/S^2$ has expectation factor at most
$1+2\mu_t/S+\nu_t/S^2$, using
$|L|/S\leq1+L^2/S^2$.  The pathwise sums just established therefore
give bounded discount loss.  The gaps from $.085a$ to $.09a$ and
from $.115a$ to $.12a$ accommodate the total-account extraction
factor.  Initial trace values satisfy the smaller early-phase bounds.

\emph{The lower bound for $\mathcal T$.}
We now control $\mathcal T-\mathcal T(l_1)$.  By
\eqref{eq:original-13}, a uniform lower comparison scale at a primary is
\[
 T_{\min}\asymp n^{-\theta/2}l_0/\alpha.
\]
This comparison uses $\alpha n/l_0\to\infty$, uniformly over the
primary points.  In \eqref{eq:original-19}, the non-$J$ terms cost
at most $C/\alpha$ per step.  For the ideal $J$ term, the relevant
row and column norms are at most $Cn^{2v_*}/\alpha^3$, its entries
are at most $Cn^{v_*}/\alpha^2$, and its full sums are centered.
Thus its cumulative absolute signed-mean bound and second-moment
bound are, respectively,
\[
 C\left(\frac{n^{2v_*}b^{3/2}}{\alpha^3l_0}
          +\frac{n^{v_*}b}{\alpha^2l_0}\right),
 \qquad
 C\left(f_b\frac{n^{4v_*}}{\alpha^6}
          +\frac{bn^{2v_*}}{\alpha^4l_0^2}\right).
\]
The nonideal $J$ coefficient is bounded using
$j_u+B_m+|D_{vv}|+|D_{uv}|$.  Its cumulative mean is at most
\[
 \frac{Cn^{v_*}}{\alpha^2}
 \left(n^\eta\sqrt b+bB_m+
                      \frac{n^\eta b}{\sqrt{l_0}}\right),
\]
and the crude cumulative second bound
$Cbn^{2v_*}/\alpha^4$ suffices.  Dividing the means by $T_{\min}$
and the seconds by $T_{\min}^2$ gives fixed power savings.  One of
the potentially tight ratios is bounded by
\[
 C\frac{b^{3/2}n^{2v_*+\theta}}
                {l_0^2\alpha_*^2}=o(1).
\]
Choose a small fixed $\kappa>0$ and apply the deviation-square
account at scale $n^{-\kappa}T_{\min}$, taking $\kappa$ smaller
than the established margins.  Its initial deviation is zero, its
discount loss is bounded, and its extracted deviation is
$o(T_{\min})$.  This proves the stated lower bound, with the
additional $n^{-\theta/2}$ cushion.  All conclusions are uniform
on their stated parameter sets: the trace-error estimates hold at
every tracked point, and the lower bound for $\mathcal T$ holds at
the primaries.
\end{proof}

\subsection{Classification of update terms}

The entry expansion used in the early phase must now be grouped for
one prescribed endpoint.  We will use this grouping both for $T_S$
and for individual entries.  For the current pair, write
\[
 \begin{gathered}
 T=U^TPU-G_0,\qquad E_{\mathrm{pair}}=U^TEU,\\
 H_P=(H_0^{-1}+T)^{-1},\qquad
 H_0=-\sigma hF+R,\quad \|R\|\leq C/l.
 \end{gathered}
\]
Thus $H_P$ is the coefficient for the numerical update of the reference
$P$, whereas $H$ uses the actual matrix $W$.  Before recentering the
output axes, the error increment is
\[
 \Delta E=-(WUHU^TW-PUH_PU^TP)-\mathcal R_P,
 \qquad \mathcal R_P=[PUH_PU^TP]_{\ge L}.
\]
The brackets select the formal terms of degree at least $L$, as in
\eqref{eq:original-1}.  We call a diagonal entry of the reference
deviation $T$ a local topology factor, and an off-diagonal entry a
cross topology factor.  Their absolute majorants are $j_u,j_v$ and
$b_{uv}$, respectively.  The entries of $E_{\mathrm{pair}}$ are
bounded by $B_m$.
On every valid pair $T$ and $E_{\mathrm{pair}}$ are small, so the
endpoint Neumann expansions converge absolutely and geometrically.

In each scalar term, the two outer entries are evaluated at selected
endpoints.  We call them \emph{two fixed outers} when both are at the
prescribed endpoint $u$, \emph{two random outers} when both are at the
averaged endpoint $v$, and \emph{distinct outers} when one is at each.
This convention concerns the endpoint indices, whether an outer entry
comes from $P$, $E$, or $W$.

\begin{lemma}[Endpoint word classification]
\label{lem:cleanup-words}
After expansion of a numerical update difference and of its formal
truncation tail, every term is covered by the following classes:
\begin{enumerate}
\item a term containing a non-pure swap $R$, and hence a factor $C/l$;
\item a pure-swap term with no intermediate endpoint factor;
\item a pure-swap term with exactly one endpoint $E$ factor and no
topology factor;
\item a term with a cross topology factor $b_{uv}$;
\item a pure-swap term with local topology factors only, and at
most one endpoint $E$ factor;
\item a term with at least two endpoint $E$ factors.
\end{enumerate}
In class~3 the single endpoint $E$ factor may be grouped with
$W$ outers.  For class~5, if both outers are at $u$, an opposite
local factor $j_v$ is necessary; the symmetric statement holds
for two outers at $v$.  If class~5 has no endpoint $E$ factor and
has distinct outers, local factors at both endpoints are necessary.
All remaining factors and choices of their positions can be summed
with a bounded geometric multiplier.
\end{lemma}

\begin{proof}
Expand
\[
 H=H_0\sum_{k\geq0}
       \bigl(-(T+E_{\mathrm{pair}})H_0\bigr)^k
\]
and the analogous reference coefficient.  Every factor is either
a pure swap, a non-pure swap, a diagonal topology entry, a cross
topology entry, or an endpoint $E$ entry.  If a non-pure swap occurs,
assign the word to class~1.  Otherwise first separate the number of
endpoint $E$ entries, and then whether any off-diagonal topology
entry occurs.  The listed classes cover all resulting cases.
The terms with no endpoint or outer $E$ that cancel numerically
leave precisely the formal tail; their endpoint classification is
unchanged by restricting total degree.

To see the useful grouping in class~3, first write the difference
of the two numerical updates as outer differences with the reference
coefficient, together with
\[
 WU H_P E_{\mathrm{pair}} H U^TW.
\]
Replacing both coefficients in this expression by their pure-swap
parts gives the asserted single-$E$ term with $W$ outers.  The
remainders contain topology, a non-pure swap, or a second endpoint
$E$, and belong to the other classes.

A pure swap changes $u$ to $v$ and $v$ to $u$.  A local topology
entry preserves its endpoint.  If both outers are at $u$, the
first and last intermediate positions of a pure-swap word are at
$v$.  With no endpoint $E$, the first such position supplies $j_v$.
With one endpoint $E$ and at least one local topology entry, there
are at least two intermediate positions, so at least one of the
first and last positions is a local topology entry at $v$.  With
distinct outers and no endpoint $E$, a nonempty local-only word
has at least two intermediate positions and visits both endpoints.
This proves the locality assertions, including the cases in which
the endpoint $E$ itself is off diagonal.

Finally the uniform small bound for each endpoint error makes the
absolute sum over additional factors geometric.  Marking one or
two factors of a specified type introduces at most a fixed power
of the word length, still summable.  The coefficientwise majorants
used for the formal tail have the same property.  Therefore all
the estimates below can be proved for the displayed factors and
then multiplied by a fixed constant.
\end{proof}

\subsection{The surviving diagonal-square statistic}

Before treating the polynomially many individual entries, we retain an
aggregate bound on their stub diagonals.  It controls averages of an
opposite-endpoint diagonal error; using only the entry supremum in those
averages would lose the required margin.

For a step with prescribed endpoint $u$, let
\[
 m_x=\mathbb E\bigl[\Delta e_x(x)
                   \mathbf1_{\{x\text{ survives}\}}\bigr]
 \quad(x\ne u),\qquad m_u=0,
\]
and let
$V_t=\mathbb E\sum_{x\text{ survives}}(\Delta e_x(x))^2$.
Expanding the surviving squares, dropping the nonnegative deleted
old squares, and applying Cauchy--Schwarz gives
\begin{equation}
 \mathbb ET_S'\leq T_S+2\sqrt{T_S}\,\|m\|_2+V_t.
\label{eq:original-37}
\end{equation}

\begin{lemma}[Cumulative diagonal-square budgets]
\label{lem:cleanup-diagonal}
Under the cleanup bootstrap, through every chosen history segment,
\[
 \sum_t\|m\|_2\leq n^\eta,\qquad
 \sum_tV_t\leq n^\eta+Cf_bT_{\mathrm{cap}},
\]
where $T_{\mathrm{cap}}$ is the chosen bootstrap cap for $T_S$.
\end{lemma}

\begin{proof}
We first estimate the vector of signed means $m$, retaining the
cancellations of the zero-intermediate and single-$E$ terms.  We then
bound $V_t$ by absolute squares.

\emph{Survival and recentering.}
We initially work on the old centered axes.  If $x$ is an eligible
partner, survival excludes the choice $v=x$.  Inserting this value
into the partner average and subsequently removing it changes its
$x$th component by at most $CB_m/l$; an ineligible $x$ needs no such
correction.  There are $O(l)$ components,
so the vector cost is $CB_m/\sqrt l$ per step.  The recentering
formula
$\xi_x'=\xi_x+(\xi_u+\xi_v)/(l-2)$ gives the same bound.
The old-axis error after a candidate update is $O(B_m)$ by the
convergent update expansion and the tail bound.  Thus these
corrections sum to $O(bB_m/\sqrt{l_0})=o(1)$.

\emph{Signed means.}
Let $I$ be the eligible partner set, and write a bar for its
uniform average.  Full sums of every centered $p$ or $e$ row
vanish.  The omitted set has size $O(b)$, so the pure-swap terms
with no intermediate factor have the following vector bounds:
\[
\begin{split}
 \|(p_x(u)\overline{e_x(v)})_x\|_2
       &\leq CbB_m/l,\\
 \|(e_x(u)\overline{e_x(v)})_x\|_2
       &\leq CbB_m/l,\\
 \|(e_x(u)\overline{p_x(v)})_x\|_2
       &\leq CB_m(1+L)^3\sqrt b/l.
\end{split}
\]
The first two use the omitted-entry bound $B_m$ and bounded column
norms.  For the third, multiply the sum of omitted $P$ columns by
the diagonal matrix with entries $e_x(u)$; its norm is at most
$B_m$, while \eqref{eq:original-7} bounds the norm of the column
sum by $C(1+L)^3\sqrt b$.

The pure $P$--$P$ zero-intermediate tail cancels over the full
partner set coefficient by coefficient.  For its omitted part,
apply the same coefficient operator estimate to the omitted
indicator vector, multiply by the bounded opposite outer entry,
and sum the polynomially many degree choices.  Its vector bound
is $n^\eta\sqrt b/l$.  Non-pure swaps contribute $C/l$ per
step: the pointwise product of two fixed columns has bounded
$\ell^2$ norm, as does the averaged product.  The cumulative
costs just listed are small, using \eqref{eq:original-35} and
$b^2B_m/l_0=o(1)$.

For the pure single endpoint $E$ term of
Lemma~\ref{lem:cleanup-words}, the two fixed outers give the
vector $(w_x(u)^2)_x$ times $\overline{e_v(v)}$.  Its norm is
bounded, and
\[
 |\overline{e_v(v)}|
      \leq C\frac{|Y_E|+\sqrt{bT_S}}l,
 \qquad Y_E=\operatorname{Tr}(\Xi^TE\Xi).
\]
The omitted diagonal sum was bounded by Cauchy--Schwarz.  For two
random outers the averaged-square vector has norm $C/\sqrt l$,
and its scalar factor is $e_u(u)$.  With distinct outers, the
average row dot against the $u$th row of $E$ is $O(1/l)$;
multiplication by the bounded-norm fixed column gives a vector
bound $C/l$.  These contributions sum to small quantities, since
\[
 |Y_E|\leq n^{.09a}/\alpha_*+Q_a,\qquad
 f_bQ_a=o(1),\qquad
 b^{3/2}\sqrt{T_{\mathrm{cap}}}/l_0=o(1).
\]

It remains to bound the absolute-mean classes.  With a cross
topology factor, two fixed outers cost $CQ_a/l$: use the
$\ell^1$ bound on the $u$th cross row and the bounded norm of
the squared outer column.  Two random outers cost $C/l$: apply
the bounded operator $(r_x(v)^2)_{xv}$ to the bounded-norm cross
row $(b_{uv})_v$, and divide by $|I|\asymp l$.  Distinct outers
also cost $C/l$, by taking the row dot at the random endpoint
before the fixed-column norm.

For local-only topology, two fixed outers require $j_v$ and
therefore cost $CQ_a/l$.  Two random outers require $j_u$ and
cost $Cj_u/\sqrt l$.  With distinct outers and no endpoint $E$,
both endpoint local factors occur; bound the random factor and
use the $\ell^1$ bound $C\sqrt l$ for its outer row, obtaining
$Cj_u/\sqrt l$.  With exactly one endpoint $E$ and local
topology, the same-outer cases still contain the opposite local
factor, while the distinct case is bounded by $CB_m/\sqrt l$.
Two endpoint $E$ factors cost $CB_m^2$ per step.  These cases
are exhaustive by Lemma~\ref{lem:cleanup-words}; outer $E$
choices are already included in $r$.  Their sums are bounded by
$n^\eta$, using Lemma~\ref{lem:cleanup-pathwise},
$bB_m^2=o(1)$, and \eqref{eq:original-35}.

\emph{Squared increments.}
For the budget $V_t$, a distinct-outer product gives
\[
 \frac{C}{l}\sum_x r_x(u)^2
                         \sum_{v\in I}r_x(v)^2\leq C/l.
\]
With two fixed outers, an opposite local or cross factor has
conditional squared mean at most $n^\eta/l$, and the fourth
sum of the fixed column is bounded.  The single opposite
diagonal $E$ factor contributes $CT_S/l$ instead.  Two endpoint
$E$ factors give $CB_m^4$.

With two random outers, a fixed local factor gives $Cj_u^2$,
using the bounded fourth-column sum for each partner; a cross
factor gives $C/l$ by its squared row bound.  Any remaining
endpoint $E$ factor can be bounded by $B_m$, giving $CB_m^2$
per step.  The no-intermediate same-end case already contains a
non-pure swap.  Thus there is no untreated same-end term without
either the stated opposite factor, the single opposite diagonal
$E$, or at least two endpoint $E$ factors.

Squares of the finite sums are bounded by the corresponding
sums of squares up to constants.  The infinite remainder is
handled by weighted Cauchy--Schwarz using its geometric decay,
as in Lemma~\ref{lem:cleanup-words}.  Centering and tail terms
obey the same majorizations.  Summing gives
$n^\eta+Cf_bT_{\mathrm{cap}}$, by
\eqref{eq:original-36}; all $bB_m^2$ and $bB_m^4$ terms are
small.  This proves the lemma.
\end{proof}

To extract the diagonal-square bound, choose a normalization $S_T$
larger by a small power than the preceding $T_S$ cap,
$(\sum\|m\|_2)^2$, and the small-loss part of $\sum V_t$.
From \eqref{eq:original-37}, the account $1+T_S/S_T$ has
expectation factor at most
\[
 1+C\|m\|_2/\sqrt{S_T}+V_t/S_T.
\]
Choose the new extraction cap $T_{\mathrm{cap}}$ a slightly
larger power than $S_T$.  Its exponent can still be prescribed
arbitrarily small; moreover $f_bT_{\mathrm{cap}}/S_T=o(1)$.
Lemma~\ref{lem:cleanup-diagonal} now bounds the accumulated
discount loss.  This establishes the required $T_S$ bound with
the usual extraction cushion.

\subsection{Individual increments and special updates}

Let $\mathcal F(P)=P-PUH_PU^TP$ denote the numerical update of
a reference matrix.  Before output recentering, the exact
difference identity is
\begin{equation}
 \mathcal F(W)-\mathcal F(P)
   =T_P(E-EUHU^TE)T_P^T,
 \qquad T_P=I-PUH_PU^T.
\label{eq:original-38}
\end{equation}
The inverse-update difference identity first gives
\[
 \mathcal F(W)-\mathcal F(P)=T_PE T_W^T,
 \qquad T_W^T=(I-UHU^TE)T_P^T,
\]
where the second identity follows by coefficient subtraction.
These are algebraic identities for
the convergent numerical updates and do not require $P$ to be
positive.  The actual error update differs by the formal tail,
whose entries are bounded by \eqref{eq:original-6}.

Call a surviving update \emph{special} for an entry if one of its
two axes has actual vertex distance at most
$R_s=\lceil\rho\log_q n\rceil$ from a selected endpoint before
the update.  Through the life of a fixed entry, at most
$Cq^{R_s}$ updates are special.  To prove this, take the actual
graph at the end of the history segment.  Every selected
endpoint responsible for a special event belongs to one of
the two radius-$R_s$ balls there.  These balls have at most
$Cq^{R_s}$ vertices, and at most $d$ selected stubs can occur
at any one vertex.  If nonspecial survival has positive
conditional probability, conditioning on it removes at most
another $Cq^{R_s}=O(b)$ possible partners.

\begin{lemma}[Level-sensitive jump bound]
\label{lem:cleanup-jumps}
For a surviving entry $E_{ij}$ of total level $k$, every allowed cleanup
step satisfies
\[
 |\Delta E_{ij}|/B_k\leq n^{-2\rho}
\]
for all sufficiently large $n$.
\end{lemma}

\begin{proof}
In the linear part of \eqref{eq:original-38}, replacing output
axis $i$ by selected stub $x$ has coefficient
\[
 -\sum_{s\in\{u,v\}}p_i(s)(H_P)_{sx}.
\]
Every transfer whose input level is at least the replaced
output level has coefficient at most $n^{-\beta_R}$ for some
fixed $\beta_R>19\rho$.  A surviving initially clean axis is
farther than $R_0$ from every selected endpoint, so both outer
entries are small by the path decay and centering estimates.
For a surviving bad axis, the only possible input of equal or
higher level is a bad first endpoint $u$.  Its outer entry at
the clean partner $v$ is small.  The remaining contribution
uses $(H_P)_{uu}$, which is small by the opposite clean
return/cross bound and the pure-swap expansion
\eqref{eq:original-21}.  Global axes have only lower-level
selected inputs.  The fixed separation margin permits the
choice $\beta_R>19\rho$ by taking $\rho$ small enough.

Consequently, apart from the persisting identity term, transfers
that all decrease level cost $CB_{k-1}$ when $k\geq1$.
Any term with a nondecreasing transfer costs
$Cn^{-\beta_R}B_m$.  The nonlinear part of
\eqref{eq:original-38} costs $CB_m^2$, and output recentering
costs $CB_m/l$.  The remaining formal tail costs $B_{\rm tail}$.
After division by $B_k$, the first bound is $Cn^{-4\rho}$,
the second is at most $Cn^{-\beta_R+16\rho}$, and the last
three are smaller than $n^{-2\rho}$ by the parameter margins.
There is no decreasing transfer when $k=0$.  This proves the
claim after absorbing fixed constants.
\end{proof}

For each sign of an entry, start its dimensionless comparator
at $1$ and increase it by $n^{-2\rho}$ at each special
surviving update.  Lemma~\ref{lem:cleanup-jumps} implies that
its signed positive excess cannot increase at such a step.
The total comparator increase is
$O(q^{R_s}n^{-2\rho})=o(1)$.

\subsection{Conditional entry budgets on ordinary updates}

The preceding jump bound pays for the few special updates.  For the
remaining updates we need a smaller cumulative signed mean and
variance.  The following estimates use the actual conditional law of
the partner, including the exclusions required for survival.

\begin{lemma}[Individual predictable entry budgets]
\label{lem:cleanup-entry-budgets}
For every fixed entry while its labels are alive, and through
every chosen history segment, there is a fixed $c_5>0$ such that
\begin{equation}
\begin{split}
 &\sum_{\substack{t:\text{nonspecial survival}\\
                         \text{has positive probability}}}
 \frac{\bigl|\mathbb E(\Delta E_{ij}\mid
                           \text{nonspecial survival})\bigr|}{B_k}
       \leq n^{-c_5},\\
 &\sum_{\substack{t:\text{nonspecial survival}\\
                         \text{has positive probability}}}
 \frac{\mathbb E((\Delta E_{ij})^2\mid
                           \text{nonspecial survival})}{B_k^2}
       \leq n^{-c_5}.
\end{split}
\label{eq:original-39}
\end{equation}
Here $k=\nu(i)+\nu(j)$ is the pair's total level.  Each sum stops when
one of its labels dies.
\end{lemma}

\begin{proof}
Conditional on nonspecial survival, the partner is uniform in
a set omitting $O(b)$ stubs.  All estimates below hold for
such sets, even when the omitted set changes with the history.
We give bounds at the smallest scale $B_c$; since $B_k\geq B_c$,
they imply the displayed assertions.

\emph{Signed means.}
For zero-intermediate pure swaps with an outer $E$, use the
centered full partner sum.  An omitted average of an $E$ row
has magnitude at most $CbB_m/l$, and an omitted average of a
$P$ row has magnitude at most $C\sqrt b/l$ by its row norm.
Lemma~\ref{lem:cleanup-pathwise} therefore gives cumulative
absolute signed-mean bounds
\[
 Cn^\eta b^{3/2}B_m/l_0,\qquad
 Cb^{3/2}B_m/l_0,\qquad
 Cb^2B_m^2/l_0
\]
for $p_i(u)\overline{e_j(v)}$,
$e_i(u)\overline{p_j(v)}$, and the two-outer-$E$ term,
respectively.  For the first term use
$\sum_t A_i^t(u_t)\leq n^\eta\sqrt b$; for the second and
third it suffices to use the entry cap on the prescribed
endpoint.  The pure $P$--$P$ zero-intermediate tail cancels
coefficientwise over all partners.  In the omitted part,
Cauchy--Schwarz gives $C\sqrt b$ for the omitted absolute
row sum, and the prescribed-endpoint path sum is
$n^\eta\sqrt b$.  Summing degree choices yields the bound
$Cn^\eta b/l_0$.

For the single endpoint $E$ term with $W$ outers, two fixed
outers contribute at most
\[
 Cn^\eta\frac{|Y_E|+\sqrt{bT_{\rm cap}}}{l_0}.
\]
Indeed their product sums over the prescribed endpoints to
at most $n^\eta$, and the signed opposite diagonal mean is
bounded as in Lemma~\ref{lem:cleanup-diagonal}.  Two random
outers cost $Cf_bB_m$, by the bounded outer row dot divided
by $l$.  Distinct outers cost $Cn^\eta\sqrt b/l_0$, taking
the random row dot first and then summing the fixed outer.

For a cross topology factor, the cumulative bounds for fixed,
random, and distinct outer endpoints are, respectively,
\[
 Cn^\eta Q_a/l_0,\qquad Cf_b,\qquad
 Cn^\eta\sqrt b/l_0.
\]
The first uses its row $\ell^1$ bound and the summed product
of the fixed outers; the second uses the bounded random
outer row dot and the bounded cross-entry supremum; the
third uses the bounded cross row norm against the random
outer before summing the fixed one.

For local-only topology, two fixed outers require the
opposite local factor and have the same $Cn^\eta Q_a/l_0$
bound.  Two random outers cost at most $Cf_b$.  With distinct
outers and no endpoint $E$, both local factors are present;
bound the random local factor by its supremum and average
its outer with cost $C/\sqrt l$.  The remaining prescribed
sum $\sum_t r_i^t(u_t)j_{u_t}^t$ is at most $n^\eta$,
so the total is $Cn^\eta/\sqrt{l_0}$.  With one endpoint
$E$ and local topology, the same-outer requirements persist;
the distinct case instead costs
$Cn^\eta B_m\sqrt{b/l_0}$.

Terms with at least two endpoint $E$ factors have the combined
bound
\[
 Cn^\eta B_m^2
                  \left(1+b/l_0+\sqrt{b/l_0}\right).
\]
For two fixed outers this follows from their pathwise product
sum; for two random outers their conditional product mean
is $O(1/l)$; for distinct outers average the random row and
sum the fixed one.  Non-pure swaps have instead the combined
bound
\[
 Cn^\eta\left(l_0^{-1}+bl_0^{-2}
                              +\sqrt b\,l_0^{-3/2}\right),
\]
by the same three outer configurations and the additional
$1/l$ factor.  Centering contributes $Cf_bB_m$.
Lemma~\ref{lem:cleanup-words} proves that these bounds cover
every mean term, including the tail.

\emph{Second moments.}
We next bound the cumulative second moments, keeping the same
three outer configurations.  With distinct outers, averaging the random
outer square gives $C/l$,
and \eqref{eq:original-36} bounds the sum of prescribed
outer squares.  The resulting budget is $Cn^\eta/l_0$.
With two fixed outers, an opposite local or cross factor
has conditional squared mean at most $n^\eta/l$; the
pathwise sum of outer product squares is at most $n^\eta$.
The resulting bound is again $Cn^\eta/l_0$, after adjusting
the preliminary loss.  The single opposite diagonal $E$
factor costs $Cn^\eta T_{\rm cap}/l_0$, while two endpoint
$E$ factors cost $Cn^\eta B_m^4$.

With two random outers, a fixed diagonal $E$ or another
otherwise unestimated endpoint $E$ contributes at most
$Cf_bB_m^2$.  A fixed local factor contributes
$Cn^\eta/l_0$, because its pathwise squared sum is bounded
and the random outer fourth-product average is $O(1/l)$.
The remaining cross topology term with two random outers requires
a joint sum over steps and partners: its outer squares must be
retained.  Separate its actual-path and uniform parts by writing
\[
 b_{uv}^t\leq b_{uv}^{t,\mathrm{act}}+U,
 \qquad U=Cn^{\eta_0}q^{L/2}/l_0,
\]
where $\eta_0$ is a sufficiently small preliminary loss and the
centering term $C/l$ is included in $U$.  The squares of these
two parts can be estimated separately, up to a fixed factor.
For the actual-path part, bound one outer by its constant supremum and
write the other squared outer as
$C(A_i^t(v)^2+B_m^2)$.  The numerator of the cumulative
partner average is thus bounded by
\[
 C\sum_v\sum_t
        \bigl(A_i^t(v)^2+B_m^2\bigr)
                         (b_{u_t v}^{t,\mathrm{act}})^2.
\]
Dominate every $A_i^t(v)$ by its terminal actual-path row plus
$U$.  This majorant has squared sum at most $n^\eta$ over the
initial stubs; its added uniform squared mass is at most
$Cl_1U^2=o(1)$.  For each fixed $v$, the sum
$\sum_t(b_{u_t v}^{t,\mathrm{act}})^2$ is at most $n^\eta$ by
Lemma~\ref{lem:cleanup-pathwise}.  Choosing the preliminary
loss smaller gives a total $n^\eta$.  The $B_m^2$ part is
at most $CbB_m^2$, using the bounded current cross-row
squared sum at every step.  Division by $l\asymp l_0$
gives $Cn^\eta/l_0+Cf_bB_m^2$.

For the uniform part, retain both outer squares.  At each
step their sum is bounded:
\[
 \sum_{v\in I}r_i^t(v)^2r_j^t(v)^2
 \leq \bigl(\sup_v r_j^t(v)^2\bigr)
                        \sum_{v\in I}r_i^t(v)^2\leq C.
\]
Its cumulative partner-average contribution is consequently at most
\[
 CbU^2/l_0=Cn^{2\eta_0}bq^L/l_0^3.
\]
This can be absorbed into the preceding $n^\eta/l_0$ budget.
Directly, its ratio to $B_c^2$ is
\[
 Cn^{2\eta_0}\frac{bq^L}{l_0^3B_c^2}
   =n^{-1/3-1.085a+2\eta_0+o(1)},
\]
which tends to zero by a fixed power.

Non-pure swaps and recentering yield smaller bounds.  As
before, geometric summability handles the longer endpoint
words, and the same-endpoint configurations without an
opposite topology factor are exactly the zero-intermediate
non-pure term, the single opposite diagonal $E$ term, or a
term with at least two endpoint $E$ factors.  This verifies
exhaustiveness of the square estimates as well.

\emph{Comparison with the target scales.}
All displayed mean budgets are smaller than $B_c$
by a fixed power, and all square budgets are smaller than
$B_c^2$ by a fixed power.  Some of the direct comparisons
that control the smaller margins are
\[
\begin{aligned}
 \frac{f_b}{B_c}&=n^{-.08a+o(1)},&
 \frac{l_0^{-1/2}}{B_c}&=n^{-.35a},\\
 \frac{l_0^{-1}}{B_c^2}&=n^{-.70a},&
 \frac{Q_a/l_0}{B_c}&=n^{-.64a+\eta+o(1)}.
\end{aligned}
\]
The bound on $Y_E$ gives another power saving, and the terms
with $B_m$, $b^{3/2}/l_0$, or $B_m^2$ have additional
margins from \eqref{eq:original-35}.  Choose all preliminary
losses, the exponent of $T_{\rm cap}$, and $\rho$ within
these strict margins.  A common positive $c_5$ then gives
\eqref{eq:original-39}.
\end{proof}

\subsection{Simultaneous extraction and completion of cleanup}

We finish the proof of Proposition~\ref{prop:cleanup}.  Give every
entry and each of its two signs its own alive-label account,
with the fixed initial family denominator.  For an entry of
total level $k$, let $c_t$ be its dimensionless comparator
defined above.  Its excess account is
\[
 1+\left(\frac{(\pm E_{ij}/B_k-c_t)_+}{\lambda}\right)^p
 \quad\text{while alive},\qquad 0\quad\text{after death},
 \qquad\lambda=n^{-\zeta}.
\]
Choose $\zeta>0$ smaller than $2\rho$, $c_5/2$, and the
individual-increment margins.  On an ordinary surviving
branch, the comparator is unchanged.  The bounded-increment
inequality \eqref{eq:original-34}, together with
\eqref{eq:original-39}, gives a pathwise summably small
discount budget.  On a special surviving branch the account
does not increase, by Lemma~\ref{lem:cleanup-jumps} and the
comparator increment.  On a death branch it drops to zero.
These statements give the unconditional account estimate;
no conditioning of the actual deterministic choice is
required.

For clarity, a step contributes to the predictable budget
whenever ordinary survival has positive conditional probability,
even if the chosen branch is special.  The
proof of \eqref{eq:original-39} bounds precisely this sum
along the chosen history by the distinct prescribed
endpoints and the current structural estimates.  Thus the
per-label discount product is bounded along that history.
Structural estimates at a chosen state come from their
separate structural accounts and do not depend on the
entry-discount estimate.  This is the order in which the
pathwise estimates enter the induction.

Initially every entry is $o(B_c)$ by the early-phase bound.
The comparator is $1+o(1)$ throughout cleanup.  There are
only polynomially many fixed entry labels and logarithmically
many spectral points.  A sufficiently large fixed even $p$
therefore makes the family-size factor after the $p$th root
smaller than the power saving supplied by $\lambda$.
The extracted excess is $o(1)$ in units of $B_k$, proving
$|E_{ij}|\leq2B_k$.

Join these accounts with the diagonal-square accounts of
Lemma~\ref{lem:cleanup-diagonal}, the scalar accounts of
Lemma~\ref{lem:cleanup-traces}, and the retained structural
accounts.  Give each normalized $T_S$, $L_0$, and $L_1$ account
weight one separately at each tracked spectral point, and each
$\mathcal T$-deviation account weight one at each primary.
There are only $O(\log n)$ such
scalar accounts.  The fixed family denominators are used for
the polynomially many individual entry labels, not for these
scalar accounts.  Thus a scalar account's extraction factor is
only the total normalized account budget $n^{o(1)}$ times its
bounded discount loss.  The phase initializations and
normalizations have the prescribed total-size cushions.
At a current state satisfying the bootstrap bounds, every available
pair has small endpoint block error,
so its next precision matrices are positive by the
two-dimensional inertia argument.  All candidate statistics
are consequently defined.  The minimizing choice preserves
the total discounted account; the preceding extraction
then supplies the bounds at the next chosen state.  This
closes the simultaneous induction and proves
Proposition~\ref{prop:cleanup}.

All choices use current data.  For each candidate pair, survival and
special status are determined by its endpoints and the current graph;
the comparator is then either unchanged or increased by $n^{-2\rho}$.
The candidate matrix entries and accounts determine their finite
conditional averages and hence their discounts before the minimizing
pair is selected.  The terminal-prefix graph is used only to prove
bounds on those accumulated discounts.  It is not data required by the
algorithm.
